\section{Experimental Protocol}\label{sec:protocol}

\textbf{Tasks and rooms.} Every method is scored on fixed task sets: 32 held-out and 32 OOD tasks for tracking, and 24 of each for navigation. Navigation episodes are flown in rooms from a test pool of 128 that is never used in training, each episode in its own room (\valNavPrivRooms\ rooms per seed).

\textbf{Adaptation.} Before adapting and after each of three adaptation stages, we fly four deterministic evaluation episodes (the policy mean) per task. A stage of gradient-based adaptation uses 10 exploratory episodes and one update, so three stages use 30 episodes. The MAML family is meta-trained for a single update, so stages 2 and 3 extrapolate beyond the trained procedure; where it matters we also report the one-update result. A PEARL stage adds one episode of context, and an RL$^2$ stage is one more episode of the trial. Non-adapting methods are evaluated once.

\textbf{Full stack.} Navigation is scored twice. Learned planners are first scored with the privileged planner, and every method is then scored on the full online stack (depth camera, map built in flight, A*) in \valNavFullRooms\ unseen rooms per seed, before and after one adaptation stage. The stage's exploration flights use the privileged planner, like a calibration flight in a known room.

\textbf{Metrics and statistics.} We report tracking RMSE in centimetres, with an early termination counting as 1.5\,m error for the rest of the episode, and navigation success: reaching within 0.5\,m of the goal before a collision or a 30\,s time-out. Learned methods are trained with \valSeedsRes\ seeds on the residual problem, \valSeedsGain\ on the gain problem and \valSeedsNav\ on navigation; intervals are 95\% normal intervals over tasks and seeds. Paired comparisons use the same tasks and seeds for both methods, with bootstrap intervals over task--seed pairs. Everything ran on one shared Intel Core i9-9980HK CPU without a GPU.
